% GENERATED FILE. Edit canonical lessons, then run npm run generate:books.
% canonical-source-hash: fnv1a64:30433a11fe8f6e47
% canonical-lessons: ES-C565-trama, ES-C565-argumento, ES-C565-capitulo, ES-C565-episodio, ES-C565-desenlace

\chapter{Capítulo uno}
\label{ch:es-capitulo-uno}
\hlchaptermodality{565}

\begin{chapteropening}
\textbf{By the end of this chapter:} \emph{I can talk about how a story is built: the plot, the storyline, the chapter, the episode and the ending.}

Complete the last lesson of chapter 565: I can talk about how a story is built: the plot, the storyline, the chapter, the episode and the ending.
\end{chapteropening}

\section[la trama]{la trama — the plot}

\label{lesson:ES-C565-trama}



\begin{quote}
Before the new one: say the Spanish for the legend, then the Spanish for the myth.
\end{quote}

\subsection*{You'll want to know: la trama}
\textbf{la trama} — \textquotedblleft{}the plot\textquotedblright{}.

\begin{quote}
La trama de la novela es muy complicada.
\end{quote}

\begin{quote}
\emph{The plot of the novel is very complicated.}
\end{quote}

\begin{cousinweb}
\emph{trama} is inherited from Latin \emph{trāma} 'the weft, the threads woven across' (EtymDB). A plot is the threads of a story woven together.
\end{cousinweb}

\begin{grammarlens}[title={What you've built}]
One more word for telling, explaining and describing at B1.
\end{grammarlens}

\subsection*{Your turn}
\begin{itemize}
\raggedright
  \item \emph{Say it:} \emph{la trama}
  \item \emph{Say it:} \emph{la trama}, once more
  \item \emph{Say it:} a sentence of your own with \emph{la trama}
  \item \emph{Recall:} say the Spanish for the legend, then the Spanish for the myth, then say \emph{la trama} again
\end{itemize}

\begin{tcolorbox}[breakable,colback=teal!4,colframe=teal!35!black,title={Before you move on}]
What does la trama mean? (\textquotedblleft{}The plot\textquotedblright{}.) Say it once more.
\end{tcolorbox}
\section[el argumento]{el argumento — the storyline, the plot summary; the argument}

\label{lesson:ES-C565-argumento}



\begin{quote}
Before the new one: say the Spanish for the myth, then the Spanish for the plot.
\end{quote}

\subsection*{You'll want to know: el argumento}
\textbf{el argumento} — \textquotedblleft{}the storyline, the plot summary; the argument\textquotedblright{}.

\begin{quote}
Me contó el argumento de la película.
\end{quote}

\begin{quote}
\emph{He told me the storyline of the film.}
\end{quote}

It also means a reason in a debate: \emph{un buen argumento}.

\begin{cousinweb}
\emph{argumento} is a learned borrowing from Latin \emph{argūmentum} 'a proof, the subject of a story' (EtymDB), from \emph{arguere} 'to make clear'.
\end{cousinweb}

\begin{grammarlens}[title={What you've built}]
One more word for telling, explaining and describing at B1.
\end{grammarlens}

\subsection*{Your turn}
\begin{itemize}
\raggedright
  \item \emph{Say it:} \emph{el argumento}
  \item \emph{Say it:} \emph{el argumento}, once more
  \item \emph{Say it:} a sentence of your own with \emph{el argumento}
  \item \emph{Recall:} say the Spanish for the myth, then the Spanish for the plot, then say \emph{el argumento} again
\end{itemize}

\begin{tcolorbox}[breakable,colback=teal!4,colframe=teal!35!black,title={Before you move on}]
What does el argumento mean? (\textquotedblleft{}The storyline, the plot summary; the argument\textquotedblright{}.) Say it once more.
\end{tcolorbox}
\section[el capítulo]{el capítulo — the chapter; the episode}

\label{lesson:ES-C565-capitulo}



\begin{quote}
Before the new one: say the Spanish for the plot, then the Spanish for the storyline.
\end{quote}

\subsection*{You'll want to know: el capítulo}
\textbf{el capítulo} — \textquotedblleft{}the chapter; the episode\textquotedblright{}.

\begin{quote}
Leí el último capítulo anoche.
\end{quote}

\begin{quote}
\emph{I read the last chapter last night.}
\end{quote}

On television it is an episode: \emph{el primer capítulo}.

\begin{cousinweb}
\emph{capítulo} is a learned borrowing from Late Latin \emph{capitulum} 'a little head, a heading' (EtymDB), from \emph{caput} 'head'.
\end{cousinweb}

\begin{grammarlens}[title={What you've built}]
One more word for telling, explaining and describing at B1.
\end{grammarlens}

\subsection*{Your turn}
\begin{itemize}
\raggedright
  \item \emph{Say it:} \emph{el capítulo}
  \item \emph{Say it:} \emph{el capítulo}, once more
  \item \emph{Say it:} a sentence of your own with \emph{el capítulo}
  \item \emph{Recall:} say the Spanish for the plot, then the Spanish for the storyline, then say \emph{el capítulo} again
\end{itemize}

\begin{tcolorbox}[breakable,colback=teal!4,colframe=teal!35!black,title={Before you move on}]
What does el capítulo mean? (\textquotedblleft{}The chapter; the episode\textquotedblright{}.) Say it once more.
\end{tcolorbox}
\section[el episodio]{el episodio — the episode}

\label{lesson:ES-C565-episodio}



\begin{quote}
Before the new one: say the Spanish for the storyline, then the Spanish for the chapter.
\end{quote}

\subsection*{You'll want to know: el episodio}
\textbf{el episodio} — \textquotedblleft{}the episode\textquotedblright{}.

\begin{quote}
Fue un episodio triste de su vida.
\end{quote}

\begin{quote}
\emph{It was a sad episode in her life.}
\end{quote}

\begin{cousinweb}
\emph{episodio} is usually traced to Greek \emph{epeisódion} 'an added part between the songs of a play', built on \emph{eísodos} 'a way in', from \emph{hodós} 'way'.
\end{cousinweb}

\begin{grammarlens}[title={What you've built}]
One more word for telling, explaining and describing at B1.
\end{grammarlens}

\subsection*{Your turn}
\begin{itemize}
\raggedright
  \item \emph{Say it:} \emph{el episodio}
  \item \emph{Say it:} \emph{el episodio}, once more
  \item \emph{Say it:} a sentence of your own with \emph{el episodio}
  \item \emph{Recall:} say the Spanish for the storyline, then the Spanish for the chapter, then say \emph{el episodio} again
\end{itemize}

\begin{tcolorbox}[breakable,colback=teal!4,colframe=teal!35!black,title={Before you move on}]
What does el episodio mean? (\textquotedblleft{}The episode\textquotedblright{}.) Say it once more.
\end{tcolorbox}
\section[el desenlace]{el desenlace — the ending, the outcome}

\label{lesson:ES-C565-desenlace}



\begin{quote}
Before the new one: say the Spanish for the chapter, then the Spanish for the episode.
\end{quote}

\subsection*{You'll want to know: el desenlace}
\textbf{el desenlace} — \textquotedblleft{}the ending, the outcome\textquotedblright{}.

\begin{quote}
El desenlace de la novela me sorprendió.
\end{quote}

\begin{quote}
\emph{The ending of the novel surprised me.}
\end{quote}

Also the outcome of a real event: \emph{un desenlace feliz}.

\begin{cousinweb}
\emph{desenlace} is built on \emph{desenlazar} 'to untie' (EtymDB), from \emph{des-} and \emph{lazo} 'a knot, a bow', which comes from Latin \emph{laqueus} 'a noose'. The ending unties the knot of the story.
\end{cousinweb}

\begin{grammarlens}[title={What you've built}]
One more word for telling, explaining and describing at B1.
\end{grammarlens}

\subsection*{Your turn}
\begin{itemize}
\raggedright
  \item \emph{Say it:} \emph{el desenlace}
  \item \emph{Say it:} \emph{el desenlace}, once more
  \item \emph{Say it:} a sentence of your own with \emph{el desenlace}
  \item \emph{Recall:} say the Spanish for the chapter, then the Spanish for the episode, then say \emph{el desenlace} again
\end{itemize}

\begin{tcolorbox}[breakable,colback=teal!4,colframe=teal!35!black,title={Before you move on}]
What does el desenlace mean? (\textquotedblleft{}The ending, the outcome\textquotedblright{}.) Say it once more.
\end{tcolorbox}
